\section{证明核查}\label{sec:proof-audit}
\paragraph{来源与状态。}
本节迁自 2026-10-13 的既有助手笔记 \path{notes/proof-audit-zh.md}，原 blob 为
\nolinkurl{6b2a63fd57b5b788360e64dca693657e6779d590}。
以下保留其数学分析及原有核查范围，不代表用户已经认可，也不代表本次格式迁移重新完成了这些核查。综合判断及其限制单列于第\ref{sec:assessment}节。

\subsection{几何实现与分次：原文背景部分}\label{audit:realization}
有限支撑归约把两个系统的简单生成元集合限制为有限集合，但不要求群有限。原文扩充 Tits 空间的退化方向，使环境对称型非退化，并保持根与余根独立；维数满足 Elias--Williamson 所使用的最小维数条件。

这一处理用于证明反射忠实性，以及证明两个不同反射的乘积满足
\[
\operatorname{im}(r_\alpha r_\beta-1)
 =\operatorname{span}(\alpha,\beta).
\]
环境型非退化不要求每个根平面上的限制非退化。不能在无限型中省略这个区别。

原文 \srclabel{bg:normalization} 的两个端点是
\[
\operatorname{grk}_A B^x=P_{x,b}(q),\qquad
\operatorname{grk}_A L_x(U_x)=q^{d(x,b)}P_{x,b}(q^{-1}).
\]
其中根的次数为一，\(U_x\) 是全部向上边，\(L_x\) 是限制映射的共同核。原文先在根次数为二的约定下跟踪 Fiebig 与 Elias--Williamson 的字符、取逆指标和移位，再代入 \(q=z^2\)。这里不能把全 costalk 与只有向上条件的核混为一谈。

既有笔记记录曾对照 Fiebig 文献\cite{Fiebig}命题 7.1(4)：stalk 的移位 \(k_i\) 与向上层的移位 \(2\ell(x)-k_i\) 对应，符合原文所用公式；并判断 Elias--Williamson\cite{EW} 的实现假设与原文构造相容。

\subsection{二面体组合：反射图重构与相对区间搬运}\label{audit:dihedral}
原文 \srclabel{dih:plane} 区分反射子群陪集的相对 Bruhat 序与环境诱导序。相对可比较推出环境可比较；反向通常不作断言。原文使用的强结论是有向反射图的识别，而不是两种偏序相等。

原文 \srclabel{bg:graph} 从覆盖边出发，使用六顶点、四层 \(1|2|2|1\) 的闭包规则恢复所有反射边。两个不同反射的乘积决定根平面，保证该规则不会产生假边；非覆盖边的三边细化与长度归纳给出反向包含。

原文 \srclabel{dih:transport} 的关键并不是任意二面体形状都自动保持。相对长度为二的端点之间恰有两条环境有向二边路径；同构保持全部反射图，因此目标相对长度也只能为二。由此逐层恢复相对覆盖和整个相对区间。既有笔记的判断是：这一步没有偷用根标签同构。

\subsection{二面体组合：较晚边计数与长度和}\label{audit:counts}
原文 \srclabel{dih:count} 证明：对于边 \(e:x\to y\)、相对顶点 \(c\ge_D y\)，在从另一系统搬运来的反射顺序下，较晚的向上边数满足
\[
h_x-h_y=k_D(e),\qquad
k_D(e)=\frac{\ell_D(y)-\ell_D(x)-1}{2}.
\]
原文 \srclabel{dih:sum} 用逆根集合之差，按包含 \(\alpha_e\) 的根平面分割其他正根，得到
\[
\sum_D k_D(e)=\frac{d(x,y)-1}{2}.
\]
这两个数值公式连接局部消失阶与全局分次。既有笔记亦参照 Dyer 的秩二手稿\cite{DyerRankTwo}；原论文重证了本次比较所需的计数和局部化结论。

既有有限检查使用原文整数模型，但独立枚举全部边、区间顶点及两个角序。对 \(I_2(m)\)、\(2\le m\le18\)，以及无限二面体绝对长度不超过 18 的范围，共检查 47,536 个实例，记录结果为全部通过。本次迁移未重新运行这些检查。检查只检验计数公式，不检验一般根系、图重构或层论结论。

\subsection{根平面局部化：模型与实际层分解}\label{audit:localization}
原文 \srclabel{loc:models} 先用二面体普通 KL 多项式均为一，识别秩一结构层模型。对有限二面体子群，可以使用标准实反射平面及固定多项式方向。

对无限二面体子群，诱导实现不一定是原始几何实现。原文没有直接套用 Elias--Williamson 的几何假设，而是调用 Fiebig\cite{Fiebig} 的 universal Coxeter system 字符定理（定理 9.3），并检查反射忠实、一般轨道及简单次正则轨道的要求。这是需要保留的外部输入，不能只写“二面体显然成立”。

原文 \srclabel{loc:relative-flabbiness} 把环境 flabbiness 传给局部化后的相对上闭集：不在根平面的边模已为零；按环境逆序补顶点时，尚未加入的较高端点不可能通过保留边连接一个已规定的较低端点。这保证延拓兼容。

原文 \srclabel{loc:decomposition} 在二维局部环 \(R=A_{(\alpha,\beta)}\) 上构造的是实际层同构
\[
B_R|_K\simeq\bigoplus_{c\in K}\mathcal O_{c,R}^{\oplus m_c},
\]
包括两个端点的限制映射，而不只是 stalk 秩相同。

下降归纳中，每个已有模型的边界像为
\(R/(\prod\alpha_e)\)，由 \(R\) 给出极小自由覆盖。给定 stalk 到边界的满射提升到这些覆盖的直和；Nakayama 保证提升满射，随后分裂，剩余自由直和项成为以当前顶点为顶的新模型。既有笔记的判断是：此论证没有要求局部化保持不可分解性，也没有假定当前底部的部分核自由。

\subsection{单边比较：余维一检测仅给包含}\label{audit:edge-bound}
原文 \srclabel{edge:bound} 设置
\[
I=\operatorname{res}_{x,e}L_x(F),\quad
J=\operatorname{res}_{y,e}L_y(G),\quad C=A/(\alpha_e).
\]
假设仅是上端点 \(L_y(G)\) 自由。其他邻接标签不与 \(\alpha_e\) 成比例，故
\[
L_y(G)\cap\alpha_e B^y=\alpha_e L_y(G),\qquad
J\simeq L_y(G)/\alpha_e L_y(G).
\]
因此 \(J\) 是自由的 \(C\)-模，且具有满泛秩。对每个根平面局部化，在同一个边模分解下比较两个像，其消失阶差正是 \(k_D(e)\)。合并得
\[
I\subseteq f_eJ,\qquad
\deg f_e=(d(x,y)-1)/2.
\]
这里用的是自由目标 \(f_eJ\) 等于其所有高度一局部化的交。没有假设 \(I\) 自反；也没有从高度一处相等推出全局相等。高余维商 \((f_eJ)/I\) 此时确实仍可能存在。

\subsection{全局比较：同时归纳与正矩阵}\label{audit:comparison}
归纳命题同时包括多项式不变性和所有搬运终段核的自由性。在比较区间 \([u,b]\) 时，所有严格较高顶点已经由更短区间的归纳假设提供自由性；底部自由性暂不使用。

取 \(0<q<1\)，令
\[
p_x(q)=q^{-d(x,b)/2}P_{x,b}(q),\quad
T=q^{-1/2}-q^{1/2}>0.
\]
逐边取消约束，归一化 Hilbert 状态从 \(p(q^{-1})\) 变为 \(p(q)\)。上一小节的包含给出每步
\[
D_{\mathrm{after}}\le (1+T E_{xy})D_{\mathrm{before}}.
\]
矩阵乘积枚举目标系统的递减反射路径，Dyer 路径公式把它识别为 \(\widetilde R'(T)\)。目标系统的互反关系为等式。较高顶点的多项式已经相同，因此在底部得到
\[
\Delta(q)\le\Delta(q^{-1}),\qquad \Delta=p_u-p'_u.
\]
交换两个系统给出反向不等式。严格次数界使 \(\Delta(q)\) 只含严格负的 \(q^{1/2}\) 次幂，而其倒数只含严格正次幂，故 \(\Delta=0\)。

端点相等后，逐步亏损向量满足
\[
0=\sum_k M_r\cdots M_{k+1}\delta_k,\qquad \delta_k\ge0.
\]
这些上三角矩阵非负且对角为一，正亏损不可能被后续矩阵消去。故各亏损为零。非零有限生成分次模的 Hilbert 级数在 \(0<q<1\) 为正，连有限长度的高余维商也不例外。因此得到 \(I=f_eJ\)，而非仅数值上的局部相等。

最后逆序加回底部约束。每次像自由于 \(A/(\alpha_e)\)，其 \(A\)-投射维数至多一。Schanuel 引理给出新核分次投射，再由分次 Nakayama 得到分次自由。至此才完成底部自由性，归纳闭合。
